\documentclass[11pt]{article}
\usepackage[a4paper,margin=26mm]{geometry}
\usepackage{amsmath,amssymb,amsthm,booktabs}
\usepackage[hidelinks]{hyperref}
\newtheorem{proposition}{Proposition}
\newcommand{\Area}{\mathcal A}
\newcommand{\PhiQ}{\Phi(Q)}
\title{Mathematical reductions for the INTLAB fence search}
\author{}
\date{September 2026}
\begin{document}
\maketitle

\section{Scope and hypotheses}
The chord, short-side, polynomial-angle, and rational-pair improvements below
are integrated into the INTLAB production classifier and schema-8 certificate
verifier. The later contractor and active-incidence ideas remain proposals.
The FLINT/Arb implementation is outside this work.
Write
\[
 D=(0,0),\quad C=(1,0),\quad B=(b_1,b_2),\quad A=(a_1,a_2).
\]
\[
 \Area=|Q|,\qquad S=2\Area=b_2+b_1a_2-a_1b_2.
\]
The relevant shapes are convex, positive-area quadrilaterals above the base,
whose longest side is \(DC\). All statements about boxes concern their
admissible subset; inadmissible points need not satisfy the geometric identities.
Let
\[
 \theta=1.0496,\qquad t=\theta^2=1.10166016,\qquad
 \eta=3t-\pi>0.
\]
The six-candidate identity is Proposition 6 of the supplied manuscript
\emph{Trapezium.pdf}. A strictly better reference competitor has already been
certified. Thus an upper bound \(\PhiQ^2\le t\) excludes maximality.

The target is fewer expensive evaluations, not artificially fewer survivors.
One may neither assume that a maximizer is a trapezium nor impose
\(a_2=b_2\), \(a_1+b_1=1\), or the final paper rectangles: those would use
the intended conclusion to prove itself. Opposite-pair equality remains
conditional on the analytic class in Proposition 17.

\section{Keep the height information in the flat-area argument}
\begin{proposition}
Put \(H=\max(a_2,b_2)\). Then
\[
 \PhiQ^2\le \frac{H^2}{\Area}.
\]
Consequently a box is excluded if interval arithmetic proves
\[
 \max(a_2^2,b_2^2)-t\Area\le0
\]
throughout its admissible subset.
\end{proposition}
\begin{proof}
Translate a vertical line continuously across \(Q\). The area on its left
varies continuously from zero to \(\Area\), so one position bisects the area.
Its intersection with \(Q\) is a segment of length at most \(H\).
After division by \(\sqrt{\Area}\) this is the claimed bound.
\end{proof}
The implemented flat-area bound follows by replacing \(H\) with
\(2\Area\), which is valid because the triangles \(DCB,DCA\) lie in \(Q\).
That replacement loses information: a box can satisfy \(H^2\le t\Area\)
even when \(\Area>t/4\).
This chord need not be the exact pair-5 construction; it is independently a
valid bisecting fence.

There is an elementary sharp box test. The polynomials
\[
 a_2^2-\frac t2(b_2+b_1a_2-a_1b_2),\qquad
 b_2^2-\frac t2(b_2+b_1a_2-a_1b_2)
\]
are convex in each coordinate separately. Their maxima over a rectangular box
are therefore attained among its 16 corners, by maximizing one coordinate
at a time. Evaluating the corner expressions with outward rounding proves
their upper bounds without derivatives or an interval quotient.
General fence functions do not enjoy this corner property.

\subsection{Use all four side directions}
For any direction, a translating line parallel to that direction has a
bisecting position. The segment length is bounded by the width of the
projection of \(Q\) onto that direction. Choose the direction perpendicular
to each side, so the projection width is the greatest distance to its
supporting line. Define the nonnegative turns
\[
 p=S-b_2=b_1a_2-a_1b_2,\qquad
 q=S-a_2=(b_1-1)a_2+(1-a_1)b_2.
\]
The four resulting heights are
\[
\begin{array}{c|c}
\text{supporting side}&\text{height}\\ \hline
DC&\max(a_2,b_2)\\
CB&\max(b_2,q)/|CB|\\
BA&\max(p,q)/|BA|\\
AD&\max(a_2,p)/|AD|
\end{array}
\]
For example, the \(CB\) test is the pair of polynomial inequalities
\[
 b_2^2-t\Area|CB|^2\le0,\qquad q^2-t\Area|CB|^2\le0.
\]
There is no need to compute a square root or divide by an interval side
length. Analogous tests apply to the other sides.

\section{A stronger short-side theorem}
The earlier short-side argument bounds the pair adjacent to a side of
length \(s\) by
\[
 F_{\rm pair}^2\le\gamma+\frac{s^2}{\Area},
\]
where \(\gamma\) is the aperture of the containing sector.
For completeness, if the sector apex triangles cut off by the short side
and its opposite have areas \(T_s,T_o\), convexity gives nesting and
\(T_s+T_o=\Area+2\min(T_s,T_o)\le\Area+2T_s\).
Writing their short-side radii as \(r_1,r_2\) gives
\[
 2\gamma T_s=\gamma r_1r_2\sin\gamma
 \le s^2(\gamma/2)\cot(\gamma/2)\le s^2.
\]
Together with the exact pair formula this proves the bound; the parallel
case is its continuous limit.

The improvement is to use the actual exclusion threshold, instead of the
fixed angle \(\pi/3\).
If \(\PhiQ^2>t\), every interior angle is \(>t\).
The sum of the two angles at the endpoints of the opposite side lies
strictly between \(2t\) and \(2\pi-2t\), so
\[
 0\le\gamma<\pi-2t.
\]
It follows that a necessary condition for \(\PhiQ^2>t\) is
\[
 \boxed{s^2>\eta\Area,\qquad \eta=3t-\pi.}
\]
This yields an inexpensive area-dependent certificate
\(s^2-\eta\Area\le0\), without first proving angle inequalities on the box.
The argument is a dichotomy: a point with a small angle is already excluded;
every other point satisfies the pair-angle estimate.

Combining this with \(\Area>t/4\), necessary by the flat-area argument, gives
the uniform exclusion
\[
 s\le \varepsilon_\ast,\qquad
 \varepsilon_\ast=\frac12\sqrt{t(3t-\pi)}
     \simeq0.212130772747.
\]
Thus \(\varepsilon=0.212\) is safe, a substantial improvement over \(0.122\).
For a rational check one may use \(\pi<355/113\):
\[
 \frac{t(3t-355/113)}4-(0.212)^2>0.000055.
\]
The old code's coarser bound \(\pi<3.142\) does not suffice for this particular
decimal choice, and its old parameter condition must be replaced, not bypassed.

Another useful unconditional formulation is
\[
 \PhiQ^2\le\frac{\pi+s^2/\Area}{3}.
\]
Indeed, if \(m\) is the smallest angle, then \(m\le\pi/2\) and
\(\gamma\le\pi-2m\). Therefore
\(\PhiQ^2\le\min(m,\pi-2m+s^2/\Area)
\le(\pi+s^2/\Area)/3\).

\section{Turn angle comparisons into low-degree polynomials}
For two rays bounding an interior angle \(\alpha\), write
\(c\) for their nonnegative cross magnitude and \(d\) for their dot product.
Since \(0<t<\pi/2\),
\[
 \alpha\le t
 \quad\Longleftrightarrow\quad
 \sin(t)d-\cos(t)c\ge0
 \quad\Longleftrightarrow\quad
 k d-c\ge0,\qquad k=\tan t.
\]
This follows from
\(c\cos t-d\sin t=|u||v|\sin(\alpha-t)\).
The constant \(k\) needs one validated evaluation at startup.
For a convex quadrilateral, the four polynomial residuals are
\[
\begin{aligned}
 E_D&=ka_1-a_2,\\
 E_C&=k(1-b_1)-b_2,\\
 E_B&=k((1-b_1)(a_1-b_1)-b_2(a_2-b_2))-q,\\
 E_A&=k(a_1^2+a_2^2-a_1b_1-a_2b_2)-p.
\end{aligned}
\]
Any one residual nonnegative throughout a box certifies a vertex bound.
These use signed turns on the admissible subset, so absolute-value
dependencies are unnecessary.

Moreover,
\[
 E_A=k(a_1^2+a_2^2)+(-kb_1+b_2)a_1+(-kb_2-b_1)a_2
\]
is affine in \(B\) and separably convex in \(A\). Its minimum over a box is
obtained by taking the four \(B\)-corners and minimizing two scalar quadratics
over the \(A\)-intervals. The minimizer of \(kx^2+\ell x\) on \([l,u]\) is
\(\max(l,\min(u,-\ell/(2k)))\).
Similarly,
\[
 E_B=k(b_1^2+b_2^2)+[-k(1+a_1)-a_2]b_1
                    +[-ka_2-1+a_1]b_2+ka_1+a_2.
\]
These formulas give sharp vertex tests without general automatic
differentiation. Interval arithmetic must enclose the clamped minimizer;
ordinary rounded minimizers alone are not proof bounds.

\section{Rational enclosures for the opposite pairs}
On an admissible shape, set
\[
\begin{array}{lll}
P_5=b_1-a_1,&C_5=a_2-b_2,&N_5=a_2p+b_2q,\\
P_6=a_1(b_1-1)+a_2b_2,&C_6=a_2+b_2-S,&N_6=a_2b_2+pq.
\end{array}
\]
The exact parallel-safe formulas simplify, when \(P_j>0\), to
\[
 F_j^2=\frac{N_j}{SP_j}K(C_j/P_j),\qquad
 K(z)=\frac{\arctan z}{z}
     =\int_0^1\frac{du}{1+z^2u^2},\quad K(0)=1.
\]
This cancels the side-length square roots and the apparent singularity at
parallel lines before interval evaluation.
For a possible above-threshold shape, both pair angles are
\(<\pi-2t<\pi/2\), so both \(P_j\) are positive. For a simple implementation
one may require their interval lower endpoints to be positive; if this is not
proved, skip the test.

The pair-angle restriction is itself an inexpensive exclusion. Put
\(\Gamma=\pi-2t\), so \(0<\Gamma<\pi/2\).
For either pair, an interval proof of
\[
 |C_j|-\tan(\Gamma)P_j\ge0
\]
proves \(\gamma_j\ge\Gamma\), hence some vertex angle is at most \(t\).
This test does not require pair equality or evaluation of an inverse tangent.
Its validity follows from the sign of \(\sin(\gamma_j-\Gamma)\) on
\(\gamma_j\in[0,\pi]\).

\begin{proposition}
For every real \(z\), with \(w=z^2\),
\[
 \frac{3}{3+w}\le K(z)\le\frac{15+4w}{15+9w}.
\]
\end{proposition}
\begin{proof}
The lower bound follows from Jensen's inequality for
\(v\mapsto(1+wv)^{-1}\) and \(\int_0^1u^2du=1/3\).
For the upper bound write
\[
 K(z)=1-w\int_0^1\frac{u^2}{1+wu^2}\,du.
\]
Under the probability measure \(3u^2du\), Jensen gives
\[
 \int_0^1\frac{u^2}{1+wu^2}\,du
 \ge\frac{1}{3(1+3w/5)}.
\]
Substitution gives the upper bound. No small-angle restriction is required.
\end{proof}
Consequently, for \(N_j\ge0,S>0,P_j>0\),
\[
 \frac{3N_jP_j}{S(3P_j^2+C_j^2)}
 \le F_j^2\le
 \frac{N_j(15P_j^2+4C_j^2)}
      {SP_j(15P_j^2+9C_j^2)}.
\]
A sufficient threshold test is the polynomial inequality
\[
 N_j(15P_j^2+4C_j^2)
   -tSP_j(15P_j^2+9C_j^2)\le0.
\]
It has no transcendental operations, square roots, or divisions.
It is only a sufficient test; unresolved boxes continue to stronger tests.
Rational bounds can be sharpened further using quadrature bounds for the
displayed integral, but this requires a separate error proof.

Pair equality is better tested jointly than by comparing two unrelated
interval values:
\[
 F_5^2-F_6^2
 =\frac{N_5P_6K(C_5/P_5)-N_6P_5K(C_6/P_6)}{SP_5P_6}.
\]
The area factor cancels from the zero test. A strictly signed numerator
excludes pair equality. Rational bounds yield sufficient polynomial sign
tests; a centered enclosure for this single residual is another option.
This differs from independently centering six complete fence values.

\section{Further geometric reductions and their logical limits}
\paragraph{Constraint propagation.}
Every possible above-threshold point satisfies all four angle lower bounds,
\(\Area>t/4\), the short-edge inequalities, convexity, and side-length bounds.
Intersecting interval consequences of these conditions can shrink a box
before subdivision. For example,
\[
 a_1<a_2/k,\qquad b_1>1-b_2/k.
\]
Using \(|AD|,|CB|\le1\) improves the global restrictions to
\[
 a_1<\cos t,\qquad b_1>1-\cos t.
\]
For a nonnegative horizontal component, the angle condition bounds it by
the side length times \(\cos t\); for a negative component the conclusion
is immediate. Contracting boxes is a proof operation: the excluded slabs
and the contracted survivor must be represented in the certificate, or the
original parent must retain a replayable contractor witness.

\paragraph{Active-side incidence.}
Lemma 12 of the manuscript states that every side of a maximizing
quadrilateral hosts an active fence endpoint. With item numbering
\(D,C,B,A,5,6\), the possible incident items are
\[
 DC:\{1,2,5\},\quad CB:\{2,3,6\},\quad
 BA:\{3,4,5\},\quad AD:\{4,1,6\}.
\]
Let \(U\) be an upper bound for the minimum of the six values throughout a
box. If every candidate incident to some side has lower bound \(>U\), no
point in the box can maximize. This is a non-optimality certificate, not a
claim that \(\Phi\le\theta\). It can be inexpensive after candidate bounds
have already been computed. It cannot be used to establish Lemma 12 itself.

\paragraph{Symmetry.}
The reflection
\((b_1,b_2,a_1,a_2)\mapsto(1-a_1,a_2,1-b_1,b_2)\)
can transfer a valid exclusion from the reflected box. Endpoints must be
rounded outward unless their exactness is established.
Adding such a fallback is sound but doubles some evaluations; it needs
benchmarking. Reflection does not justify silently dropping the corresponding
half of an already recorded full-root proof.

\paragraph{Do not assume two active vertex angles.}
The discussion preceding Proposition 17 (manuscript p.~16) explicitly leaves
two cases to the certified computation: P2V1 with $V=\{C\}$ and P2V2 with
$V=\{C,D\}$. In Section 4, the P2V1/$\{C\}$ argument invokes Proposition 17
in Step 6(ii) (pp.~29--30). Its exclusion is therefore not an independent
analytic premise for the computation proving that proposition. Imposing two
active angles in that computation would be circular. Section 5's analytic
exclusion of P2V1 applies to trapezia, not all quadrilaterals under consideration.

The implemented P2V0 rule below excludes no-active-vertex configurations
without assuming pair equality. The remaining single-vertex branch could
additionally be tested against its
independently derived stationarity identity (50). The two-vertex condition
can be imposed on the other branch only. These latter additions are not enabled.

\paragraph{Implemented exclusion of all vertex fences being inactive.}
Suppose an interval upper bound $U$ for one squared opposite-pair value
$F_j^2$, $j\in\{5,6\}$, satisfies
\[
       \alpha_v>U\qquad(v=A,B,C,D)
\]
throughout the admissible subset of a box. Since $\Phi(Q)^2\le F_j^2\le U$,
no vertex fence is active. By Proposition 6, only opposite-pair fences can
be active. If only one pair is active, two sides host no active endpoint,
contrary to Lemma 12. If both are active, Sections 4 and 5 independently
exclude P2V0 for non-trapezia and trapezia with exactly one pair of parallel
sides, respectively. Remark 19 independently excludes parallelograms.
Thus no admissible shape in the box is a maximizer. This argument neither
invokes Proposition 17 nor assumes that the two pair values are equal.
It is a non-optimality proof, not necessarily a bound $\Phi(Q)\le\theta$.

The implementation obtains $U$ from the rational pair upper bounds proved
above, requiring positive $S,P_j$ and nonnegative $N_j$ throughout the box.
Invalid or nonfinite pair bounds are skipped. If $0\le U<\pi/2$, put
$k=\tan U$ with validated rounding. The strict tests are
\[
           E_v(k)=k d_v-c_v<0\qquad(v=A,B,C,D),
\]
where $d_v,c_v$ are the signed dot and cross expressions above.
They imply $\alpha_v>U$ because $\alpha_v\in[0,\pi]$ and $U<\pi/2$.
For $D,C$ these are linear interval tests. For $A,B$ the residuals are
separately convex in all four coordinates, so their maxima occur among the
16 corners. Corner expressions are evaluated with outward rounding and
every nonfinite value fails closed. A nonnegative upper bound is inconclusive;
strict negativity is required for every vertex.

The certificate status is \texttt{p2v0\_nonoptimal}, the item selects $j=5$ or
$6$, and the verifier reconstructs the bound and all four strict tests.
Schema 7 records the independent enabling flag \texttt{p2v0\_cert}; schemas
4--6 do not enable this rule. Priority is after admissibility, short-side and
flat-area tests, but before other geometric or exact/centered fence work.
The classifier attempts the rule only at normalized width at most $1/32$:
coarse tests gave no benefit. This binary64 scheduling decision only skips
an optional proof attempt and has no effect on the validity of an exclusion.

\paragraph{Implemented common active-pair/vertex value.}
Use the independent analytic reductions of P0 and P1 in Sections 4--5
(as applicable to non-trapezia and trapezia), the P2V0 exclusions there,
and Remark 19 for parallelograms.
After these reductions a maximizer must have both opposite pairs and at
least one vertex active. This is stronger than equality of the pair values
alone: their common value must be the minimum fence value. It does not
assume Proposition 17 or the P2V1 conclusion that invokes that proposition.

Let $I_1,\ldots,I_6$ enclose the normalized fence \emph{lengths} on a cell,
with opposite pairs indexed by 5 and 6, and put $J=I_5\cap I_6$.
Any maximizer in the cell must satisfy, for some $v\in\{1,2,3,4\}$,
\[
           \Phi(Q)=F_5(Q)=F_6(Q)=F_v(Q)\in J\cap I_v.
\]
Consequently the cell contains no maximizer if
\[
 J=\varnothing\quad\text{or}\quad
 \bigl(J\cap I_v=\varnothing\ \text{for every }v=1,2,3,4\bigr).
\]
For nonempty $J=[\ell,u]$, each vertex must satisfy either
$\sup I_v<\ell$ or $\inf I_v>u$ to certify this exclusion.
In particular, vertices below $J$ need not surpass it to be incompatible.
The comparisons are strict: touching closed endpoints is inconclusive.
All entries must use the same normalization and units; raw vertex angles
cannot be compared with pair lengths. Nonfinite bounds fail closed, except
that an empty pair intersection needs no vertex bound. Taking maxima and
minima of already outward-rounded endpoints introduces no new arithmetic
rounding error. This proves non-optimality, not necessarily $\Phi\le\theta$.

The implementation reuses pair intervals, lazily evaluates natural vertex
lengths before derivatives, and stops at the first possible intersection.
If centering tightens a bound, it retries with the tightened enclosures.
The optional test adds no gradient objects and never contracts a cell.
The older polynomial P2V0 prefilter remains a separate inexpensive rule.
The new status is \texttt{active\_pair\_vertex\_incompatible}, with item 0;
the stored intersection is diagnostic only. Schema 8 records an independent
enabling flag and a fixed analytic-reduction identifier. Verification
reconstructs the proof from the cell; schemas 4--7 disable the new rule.
At least \emph{one} active vertex is required, never two.

\paragraph{Polynomial range bounds.}
The low-degree residuals above are better targets for monotonicity tests,
explicit quadratic minimization, or Bernstein coefficient bounds than the
full six-fence expressions. Bernstein coefficient enclosures provide range
bounds, as illustrated in the official
\href{https://www.tuhh.de/ti3/rump/intlab/demos/html/dintlab.html}
{INTLAB demonstration}. Multivariate range bounds must use outward-rounded
coefficient transformations; checking only polynomial corner values is
valid only with a proved property such as the separate convexity above.

\section{Historical diagnostics and suggested order}
The diagnostic scripts and original measurement README described in this
section are preserved locally under the git-ignored \texttt{Intlab/Archive/}.
They are not part of the production driver. The current verified centered
chain, hulls, and stage timings are documented in \texttt{Intlab/README.md}.
The historical \texttt{assess\_exclusions.m} reads every saved survivor
from a specified certificate and evaluates the proposed tests with INTLAB.
It neither modifies the ongoing search nor writes a proof certificate.
It reports individual counts and their union; overlapping counts must not
be added. In particular these are the complete base survivor set, not the
clustered sample of 48 subtrees discussed earlier in the thread.

On all 77,910 saved base survivors the sharp height test excludes 12,617
cells, and the four side-direction chord tests together exclude 13,102.
The improved uniform short-side test excludes 149, its area-dependent version
847, the sharpened vertex tests 58, and the rational pair tests 154 and
1,524 respectively. The pair-angle test excludes 26. The union of the
unconditional exclusions is 15,209 cells (19.52 percent), leaving 62,701.
The 15 pair-disjointness exclusions are already in this union.
The archived theory README records the input hash and the complete table.

The most promising implementation order is: height and side-direction chord
tests; the strengthened short-side rule; polynomial vertex tests; rational
pair tests and the common-denominator equality residual; then contractors.
The first four groups are implemented in the existing leaf-classification
architecture, with the common area cancelled for rational pair equality.
Contractors require a more substantial certificate extension.
Counts on saved cells do not establish an end-to-end runtime improvement,
and none of these changes justifies a new speed claim without a full
comparison on the same input seeds.

A bounded serial implementation comparison on 64 evenly spaced saved base
survivors, refined to normalized width $0.0078125$, gave 467 leaves and
163 survivors under the old rules, versus 391 leaves and 114 survivors under
the new rules. Search times were approximately 14.5 and 15.8 seconds.
The extra exclusions therefore did not yet yield a speedup on this step;
their possible benefit over a longer refinement chain remains to be measured.
The implemented cascade places the linear vertex and sharp height tests
before exact pairs, and the remaining algebraic tests afterward. It avoids
duplicate natural vertex evaluations and shares interval work within rule
families. No contractor or final trapezium identity is imposed.
\end{document}
